\subsection{Low-Dimensional State Sufficiency}

After local T1 survival was established, we tested whether compact state
variables were sufficient to account for first- or second-moment structure in
the local tangential non-affine residual. The grouped validation unit was the
observation: models were trained on all but one observation and evaluated on
the held-out observation. The baseline was a radius-only binned model. The
state model added $C(t)$ and $\dot{C}(t)$ to the radius covariate. Performance
was measured by incremental
\begin{equation}
R^2_{\mathrm{inc}}
=
1-
\frac{\mathrm{MSE}_{\mathrm{state}}}{\mathrm{MSE}_{\mathrm{radius}}}.
\end{equation}
A within-observation circular shift of $C$ and $\dot{C}$ provided a temporal
state-null comparison.

The radius variable depended on the unit of analysis. In the vector-level
moment test, the radius covariate was the focal individual's swarm-centered
radius. In the state-matched event-locality and history tests described below,
the radius coordinate was the swarm-level robust-z $R_{\mathrm{rms}}(t)$.

\subsection{State-Matched Event-Locality}

We next asked whether true transition timestamps carried information beyond
continuous compact state. For each true event, near-pre T1 activity over
$[-0.25,0.00]$ s was compared with same-observation non-event frames matched
in robust standardized $(C,\dot{C},R)$ space. Candidate controls were ranked
by Euclidean distance in this standardized state space, and up to five nearest
matched controls were used for each event. Candidate controls were required to
be at least 0.75 s from true transition times, and the best match distance had
to be no greater than 0.75 for the event to be accepted. Controls were not
removed after use, so the same non-event frame could in principle match more
than one event. A shifted-event null with 80 within-observation replicates was
used to test whether true event timings exceeded a temporal reference.
This endpoint-inclusive state-matched window follows the 4100 implementation
and is distinct from the half-open phase-bin convention used for event-aligned
profiles, where the transition frame was assigned to the near-post bin.

\subsection{Recent-History and Observation Heterogeneity}

Finally, we tested whether recent state-path direction supplied information
not contained in the current state. Frames were paired within the same
observation when their current $(C,\dot{C},R)$ states were close but their
recent paths through the $C$--$R$ plane were contrasting. For history window
$h$, the recent path vector was
\begin{equation}
\Delta X_h(t)=
\left[C(t)-C(t-h),\; R(t)-R(t-h)\right],
\end{equation}
and the contrast angle was the angle between two such vectors. The primary
history window was 0.50 s, the current-state distance threshold was 0.50, the
history-angle contrast threshold was 90 degrees, and paired frames were
required to be at least 1.0 s apart. The observed history-conditioned
separation was compared with within-observation shuffled-history controls.

All cross-observation summaries preserved observation identity. Observation
heterogeneity was treated as part of the empirical result, while metadata
associations were used descriptively only.
